\subsection{Padding and Radius}
\label{sec:padding}

\subsubsection{Block Padding}

Vertical (block-axis) padding follows directly from the height formula:
\begin{equation}
p_{\text{block}} = d \cdot w \cdot U.
\end{equation}

At $w = 0$, block padding is zero. At $w = 1$ with default density, $p_{\text{block}} = 1.5\,U = 6\,\text{px}$.

\subsubsection{Inline Padding}

Horizontal (inline-axis) padding uses a width multiplier derived from the wrapping level:
\begin{equation}
p_{\text{inline}} =
\begin{cases}
0 & w = 0, \text{ unbounded inline} \\
2 \cdot d \cdot U & w = 0, \text{ bounded inline} \\
\lceil 3 / w \rceil \cdot d \cdot w \cdot U & w \ge 1
\end{cases}
\end{equation}

The $\lceil 3 / w \rceil$ factor produces the following multipliers:

\begin{center}
\begin{tabular}{ccc}
\toprule
$w$ & $\lceil 3/w \rceil$ & $p_{\text{inline}}$ at $d = 1.5$ \\
\midrule
1 & 3 & $4.5\,U$ \\
2 & 2 & $6\,U$ \\
3 & 1 & $4.5\,U$ \\
\bottomrule
\end{tabular}
\end{center}

This non-monotonic behavior is intentional: single-line controls ($w = 1$) and structural sections ($w = 3$) share the same inline padding, while multi-line blocks ($w = 2$) receive wider padding to accommodate wider content.

For bounded inline elements at $w = 0$ (e.g., tag, badge, inline code), inline padding is $2dU = 3\,U$ at default density. Unbounded inline content (plain text, icon) receives no padding.

\subsubsection{Border Radius}

Border radius is defined to equal the block padding:
\begin{equation}
r = p_{\text{block}} = d \cdot w \cdot U.
\end{equation}

This identity ensures that the corner curvature is always proportional to the surrounding whitespace, maintaining consistent visual weight across density levels.

\subsubsection{Spacing Reference}

Table~\ref{tab:spacing-reference} shows all spatial properties at default density.

\begin{table}[ht]
\centering
\caption{Spatial properties at default density $d = 1.5$, $n = 1$. Values in $U$ (px at $16\,\text{px}$ font size).}
\label{tab:spacing-reference}
\begin{tabular}{lcccc}
\toprule
Property & $w = 0$ & $w = 1$ & $w = 2$ & $w = 3$ \\
\midrule
Height        & 6 (24)   & 9 (36)   & 12 (48)  & 15 (60)  \\
$p_{\text{block}}$   & 0 (0)    & 1.5 (6)  & 3 (12)   & 4.5 (18) \\
$p_{\text{inline}}$  & 3 (12)*  & 4.5 (18) & 6 (24)   & 4.5 (18) \\
Radius        & 0 (0)    & 1.5 (6)  & 3 (12)   & 4.5 (18) \\
\bottomrule
\end{tabular}

\vspace{0.5em}
\footnotesize * Bounded inline only (tag, badge, code). Unbounded inline has zero padding.
\end{table}

